% Generated draft: requires mathematical and source review.
\documentclass[11pt]{article}
\usepackage[margin=25mm]{geometry}
\usepackage{fontspec}
\setmainfont{Latin Modern Roman}
\usepackage{amsmath,amssymb,mathtools,unicode-math}
\setmathfont{Latin Modern Math}
\usepackage{enumitem,needspace,xurl,hyperref}
\setlength{\parindent}{0pt}
\setlength{\parskip}{5pt}
\setlength{\emergencystretch}{3em}
\newcommand{\N}{\mathbb{N}}
\newcommand{\Z}{\mathbb{Z}}
\newcommand{\Q}{\mathbb{Q}}
\newcommand{\R}{\mathbb{R}}
\newcommand{\C}{\mathbb{C}}
\newcommand{\F}{\mathbb{F}}
\newcommand{\A}{\mathbb{A}}
\newcommand{\PP}{\mathbb{P}}
\newcommand{\E}{\mathbb{E}}
\newcommand{\eps}{\varepsilon}
\newcommand{\dd}{\,\mathrm{d}}
\begin{document}
\section*{UnsolvedMath ID: 30000560}
\textbf{Source problem code:} catalogue entry
\subsection{Definitions and mathematical statement}
Let $\mathbb T^4=\R^4/\Z^4$, with normalized Lebesgue probability measure $\lambda$. Let $\alpha_0$ be the linear torus action associated in the source to a real-diagonalizable maximal abelian subgroup of $\operatorname{Sp}(4,\Z)$. Here $\operatorname{Sp}(4,\Z)$ preserves the standard nondegenerate alternating form on $\Z^4$. The identification or finite-index parametrization of that subgroup by $\Z^2$ must be specified; the draft uses the source's notation $\alpha_0(n)$ for $n\in\Z^2$ without inventing it.

Let $\alpha$ be a $\Z^2$-action on $\mathbb T^4$ such that $\alpha(n)$ is homotopic to $\alpha_0(n)$ for every $n\in\Z^2$. Assume the semiconjugacy $h$ in the source is given, with
\[
 h\circ\alpha(n)=\alpha_0(n)\circ h\qquad(n\in\Z^2).
\]
Under the standard convention $h$ is a continuous surjection; its existence is part of the supplied setup, not deduced here solely from homotopy. Let $\mu$ be an $\alpha$-invariant Borel probability measure such that
\[
 h_*\mu=\lambda.
\]
The requested conclusion has both parts:
\[
 \mu\ll\lambda,
 \qquad
 \mu=\omega\wedge\omega
\]
for a Lebesgue-measurable exterior two-form $\omega$ that is closed in the appropriate sense intended by the report.

Absolute continuity means that every Lebesgue-null Borel set is $\mu$-null, equivalently $\mu=f\lambda$ for a nonnegative integrable density with integral one. The form equality should mean equality of measures: $\mu(A)=\int_A\omega\wedge\omega$ for Borel sets, with the resulting top-degree form integrable and nonnegative for the chosen orientation. The source's normalization is exactly $\omega\wedge\omega$; no factor $1/2$ is inserted.

The differentiability and dynamical hypotheses on $\alpha$ are missing, as are possible irreducibility or hyperbolicity requirements on the linear action. The meaning of closedness for a merely measurable form is also unspecified. Distributional closedness $\dd\omega=0$ would require an appropriate integrability/current convention, not supplied here. No smoothness, nondegeneracy, invariance of $\omega$, ergodicity of $\mu$, or uniqueness of $\mu$ is added. The fact that the linear model lies in a symplectic matrix group does not alone state that the nonlinear action preserves a given smooth symplectic form.
\subsection{Short English statement}
For an invariant measure projecting to Lebesgue measure under the given torus semiconjugacy, the source asks for absolute continuity and a representation as the square of a suitably closed measurable two-form. The action regularity and weak closedness conventions remain essential missing hypotheses.
\subsection{Sources}
\begin{itemize}
\item Local UnsolvedMath registry, permanent ID 30000560, code catalogue entry, ``Absolute Continuity for Symplectic $\Z^2$-Actions.'' Proposed year 2006; proposer unspecified.
\item \emph{Geometric Group Theory, Hyperbolic Dynamics and Symplectic Geometry} (2007), Oberwolfach Reports, pp.\ 1991--2058, DOI 10.4171/owr/2006/33; \url{https://doi.org/10.4171/owr/2006/33}. The displayed report year 2007 differs from the DOI and proposed year 2006. This is the sole supplied literature reference.
\item Statement status corrected\_verified; the supplied verification says the official report was checked and the torus, linear action, homotopy, semiconjugacy, and desired measure form restored.
\item Literature triage checked 2026-08-21: open. It generically describes the item as an unresolved problem/conjecture/direction and says no authoritative resolution was found by the supplied targeted check. It supplies no further regularity hypotheses. Created 2026-08-25; updated 2026-09-16. Current status and source claims require independent verification.
\item Registry curation by Ulam AI / UnsolvedMath Contributors, CC BY 4.0; the underlying report retains its supplied provenance. Unreviewed statement-only draft.
\end{itemize}
\end{document}
